% ---
% id: meat-core/lem:estimate
% kind: lem
% title: "Upper estimate of the value at any feasible point"
% summary: >-
%   The objective at any exactly feasible point, a state or the Choi operators
%   of [[thm:choi]] that parametrise one, is at least the value h_low of the
%   program. It warrants the estimate every run reports, a
%   solver diagnostic and never the lower bound a key length needs.
% prominence: standard
% status: implemented
% uses: [def:feasible, thm:choi, def:admissible-declaration]
% mentions: [def:admissible-acceptance-set]
% assumes: []
% refs: []
% provenance: newly-derived
% implements: certiqs_key.meat_core.objective.objective:objective_value
% lean:
% notes: "What a run must check for the warrant to hold, membership of its point, and at which window, is the solver submodule's and is not stated here."
% ---

The number a run reports as its \tref{estimate} is the objective at the
point its optimizer returned. What that number is worth is one line, and it
is a property of the program, not of any method that searches it. The
membership $x\in\acSet\times\Feas^{J}$ is exactly what must be checked for
this warrant to hold, and which check establishes it is the solver
submodule's claim, not the core's. Definition~\ref{def:admissible-acceptance-set} is not needed:
nothing in the lemma reads convexity or closedness, and at an empty
$\acSet$ the statement is vacuous. In either form the number is an \emph{upper}
estimate of $\hlow$ and never a bound on it: nothing here says how far above
$\hlow$ it lies.

\begin{lemma}\label{lem:estimate}
Let $\Decl$ satisfy Definition~\ref{def:admissible-declaration}, let $\acSet$ be any set
of probability vectors on the acceptance alphabet, and let
$\alpha\in(1,\infty)$.
\begin{enumerate}[label=(\roman*),leftmargin=2.2em]
\item \emph{(State form.)} For every $(\mathbf q,\omega)\in\acSet\times\Feas$,
\begin{equation}\label{eq:estimate}
  O_{\alpha}(\mathbf q,\omega)\;\ge\;\hlow .
\end{equation}
\item \emph{(Choi form.)} Write
$O^{J}_{\alpha}(\mathbf q,J)\coloneqq O_{\alpha}(\mathbf q,\chi_{\Qsys}(J))$
for the objective of the Choi program \eqref{eq:choi-program} of
Theorem~\ref{thm:choi}, with $\chi_{\Qsys}$ the map of that theorem. For
every $x=(\mathbf q,J)\in\acSet\times\Feas^{J}$,
\begin{equation}\label{eq:estimate-choi}
  O^{J}_{\alpha}(x)\;\ge\;\hlow .
\end{equation}
\end{enumerate}
\end{lemma}

\begin{proof}
(i) $\hlow$ is by \eqref{eq:hlow-def} the infimum of $O_{\alpha}$ over
$\acSet\times\Feas$, and $(\mathbf q,\omega)$ is a point of that set; an
infimum is a lower bound on every value it is taken over.

(ii) By Theorem~\ref{thm:choi}, $\chi_{\Qsys}$ maps $\Feas^{J}$ onto
$\Feas$ and the value of the Choi program is $\hlow$; so for
$x=(\mathbf q,J)\in\acSet\times\Feas^{J}$ the pair
$(\mathbf q,\chi_{\Qsys}(J))$ lies in $\acSet\times\Feas$ and
$O^{J}_{\alpha}(x)=O_{\alpha}(\mathbf q,\chi_{\Qsys}(J))\ge\hlow$ by (i).
\end{proof}
